\documentclass[11pt]{article}
\usepackage[margin=1in]{geometry}
\usepackage{amsmath,amssymb,amsthm,graphicx,microtype,enumitem}
\usepackage[hidelinks]{hyperref}

\newtheorem{theorem}{Theorem}
\newtheorem{lemma}[theorem]{Lemma}
\theoremstyle{remark}
\newtheorem{remark}[theorem]{Remark}
\newcommand{\T}{\mathbb T}
\newcommand{\R}{\mathbb R}
\newcommand{\norm}[1]{\left\lVert#1\right\rVert}

\title{A continuous uniform-BV kernel whose Mercer singular-value expansion is not absolutely convergent}
\author{Candidate full counterexample to the belief stated in arXiv:2409.16453}
\date{13 August 2026}

\begin{document}
\maketitle

\begin{center}
\fbox{\parbox{0.91\textwidth}{\textbf{Status.}
Candidate full counterexample, likely valid.  There is a real continuous
kernel on $[-\pi,\pi]^2$ whose sections have uniformly bounded variation and
whose Mercer singular-value expansion converges uniformly, but not absolutely
at $(\pi/4,\pi/4)$.  Its positive singular values are simple, so the failure
cannot be removed by choosing a different singular basis.}}
\end{center}

\section{The source statement}

Jeong and Townsend consider a continuous kernel $K:[a,b]\times[c,d]\to\R$
of \emph{uniform bounded variation}, meaning that its weak partial derivatives
satisfy
\[
 \sup_y\int_a^b |\partial_xK(x,y)|\,dx<\infty,
 \qquad
 \sup_x\int_c^d |\partial_yK(x,y)|\,dy<\infty.
\]
They prove almost-everywhere and almost-uniform convergence results and then
state their belief that the Mercer singular-value expansion of every such
kernel converges ``pointwise, absolutely, and uniformly''
\cite[p.~2]{JeongTownsend}.  The complete context, including their definition,
is reproduced in Figure~\ref{fig:source}.

\begin{figure}[p]
\centering
\includegraphics[width=\textwidth]{figures/open_problem_crop.png}
\caption{Full-width rendering of PDF page 2 of arXiv:2409.16453v2.  The
uniform-BV definition is equation (3), and the stated belief appears in the
paragraph immediately below it.}
\label{fig:source}
\end{figure}

\section{Proof intuition}

Absolute SVE convergence would force the values of the singular functions to
compensate for a borderline nonsummable singular spectrum.  Trigonometric
singular functions allow us to defeat that compensation at one explicit
point.  We choose singular coefficients
\[
 a_n=\frac{1}{n\log(n+2)},
\]
which are square-summable but not summable on the odd integers.  The resulting
sine series nevertheless converges uniformly because $n a_n\to0$.

The only subtlety is bounded variation: formally differentiating loses the
factor $1/n$ and produces coefficients $b_n=1/\log(n+2)$, which are not
summable.  Convexity of $(b_n)$ resolves this.  A positive mixture of the
Fej\'er kernels has exactly these Fourier coefficients and finite $L^1$ norm.
Its mean-zero part is therefore the weak derivative of a continuous periodic
primitive $f$.  The identity
\[
 K(x,y)=\frac{f(x+y)+f(x-y)}2
\]
then supplies the uniform variation bounds, while orthogonality identifies
the series as the actual SVE.  At $(\pi/4,\pi/4)$ its absolute values contain
one half of the divergent odd subseries of $\sum a_n$.

\section{The counterexample}

\begin{theorem}\label{thm:main}
There is a real continuous asymmetric kernel
$K:[-\pi,\pi]^2\to\R$ of uniform bounded variation such that its Mercer
singular-value expansion converges uniformly on the square but is not
absolutely convergent at $(\pi/4,\pi/4)$.  Its nonzero singular values are all
simple.
\end{theorem}

We prove the theorem in four steps.

\subsection{An integrable function with prescribed slow Fourier coefficients}

Put
\[
 b_n=\frac1{\log(n+2)}\quad(n\geq0),
 \qquad
 \Delta^2b_N=b_N-2b_{N+1}+b_{N+2}.
\]
The function $x\mapsto1/\log(x+2)$ is strictly decreasing and
\[
 \frac{d^2}{dx^2}\frac1{\log(x+2)}
 =\frac{\log(x+2)+2}{(x+2)^2\log^3(x+2)}>0.
\]
Thus $b_n\downarrow0$ and $\Delta^2b_N>0$.  Let
\[
 F_N(t)=1+2\sum_{k=1}^N\left(1-\frac{k}{N+1}\right)\cos(kt)
 =\frac1{N+1}\left(\frac{\sin((N+1)t/2)}{\sin(t/2)}\right)^2
\]
with the continuous interpretation at multiples of $2\pi$.  These are the
nonnegative Fej\'er kernels \cite{Fejer}, normalized by
$\int_{-\pi}^{\pi}F_N=2\pi$.  Define, in $L^1(\T)$,
\begin{equation}\label{eq:g}
 g(t)=\frac12\sum_{N=0}^{\infty}(N+1)\Delta^2b_N F_N(t).
\end{equation}

To justify convergence and compute the coefficients, write
$d_N=b_N-b_{N+1}$.  For finite $M$,
\begin{align}
 \sum_{N=0}^M(N+1)\Delta^2b_N
 &=b_0-b_{M+1}-(M+1)d_{M+1},                                      \label{eq:mass}\\
 \sum_{N=n}^M(N+1-n)\Delta^2b_N
 &=b_n-b_{M+1}-(M+1-n)d_{M+1}.                                   \label{eq:coeff}
\end{align}
The mean-value theorem gives $N d_N\to0$.  Hence the nonnegative terms in
\eqref{eq:g} have total $L^1$ norm $\pi b_0$, and the series converges in
$L^1$.  Equations \eqref{eq:mass}--\eqref{eq:coeff} show that $g$ has constant
term $b_0/2$ and cosine coefficient $b_n$ for every $n\geq1$.  Consequently
\begin{equation}\label{eq:h}
 h=g-\frac{b_0}{2}
\end{equation}
is even, belongs to $L^1(\T)$, has mean zero, and has cosine coefficients
$b_n$.

\subsection{A continuous periodic primitive and its uniform series}

Define $f(t)=\int_0^t h(s)\,ds$ on $[-\pi,\pi]$ and extend it periodically.
Evenness and zero mean of $h$ give $f(-t)=-f(t)$ and
$f(-\pi)=f(\pi)=0$.  Thus $f$ is continuous, periodic, and absolutely
continuous, with weak derivative $h$.  Integration by parts shows that its
$n$th sine coefficient is
\begin{equation}\label{eq:an}
 a_n=\frac{b_n}{n}=\frac1{n\log(n+2)}.
\end{equation}

For completeness, we verify uniform convergence rather than cite a criterion.
If $(a_n)$ is positive decreasing and
$A_N=\sup_{n\geq N}n a_n\to0$, then for $0<t\leq\pi/2$ split a tail at
$q=\lfloor1/t\rfloor$.  On $N\leq n\leq q$, the bound
$|\sin(nt)|\leq nt$ gives a total at most $A_N$.  On the remaining indices,
Abel summation and
\[
 \left|\sum_{n=r}^s\sin(nt)\right|\leq\frac{2}{|\sin(t/2)|}
\]
give at most a fixed multiple of $A_N$.  If $t>1/N$, the same Abel estimate
applied from $N$ gives the bound directly.  For $\pi/2\leq t\leq\pi$, the
partial sine sums are uniformly bounded, so Abel summation again applies.
The endpoints are trivial.  Therefore
\begin{equation}\label{eq:uniformtail}
 \sup_{t\in\R}\left|\sum_{n=N}^M a_n\sin(nt)\right|\leq C A_N
 \qquad(M\geq N)
\end{equation}
for an absolute $C$.

Here $n a_n=b_n\to0$, so $\sum_{n\geq1}a_n\sin(nt)$ converges uniformly to a
continuous periodic function $S$.  Both $S$ and $f$ are odd and have the same
sine coefficients.  Their difference has every Fourier coefficient zero.
Convolving that continuous difference with $F_N/(2\pi)$ gives zero for every
$N$, while the displayed positive-kernel formula makes these convolutions
converge uniformly to the difference.  Hence
\begin{equation}\label{eq:fseries}
 f(t)=\sum_{n=1}^{\infty}a_n\sin(nt)
\end{equation}
uniformly on $\R$.

\subsection{The kernel and its uniform variation bounds}

Set
\begin{equation}\label{eq:K}
 K(x,y)=\frac{f(x+y)+f(x-y)}2
 =\sum_{n=1}^{\infty}a_n\sin(nx)\cos(ny).
\end{equation}
The identity and uniform convergence follow from \eqref{eq:fseries}.  In
particular $K$ is real and continuous.  Its weak partial derivatives are
\[
 \partial_xK(x,y)=\frac{h(x+y)+h(x-y)}2,
 \qquad
 \partial_yK(x,y)=\frac{h(x+y)-h(x-y)}2.
\]
Because an interval of length $2\pi$ contains one full period of every
translate of $h$,
\begin{align*}
 \int_{-\pi}^{\pi}|\partial_xK(x,y)|\,dx&\leq\norm{h}_{L^1(\T)},\
 \int_{-\pi}^{\pi}|\partial_yK(x,y)|\,dy&\leq\norm{h}_{L^1(\T)}.
\end{align*}
The bounds are uniform in the fixed variable, exactly as required by the
source definition.  The kernel is asymmetric: $K(0,y)=0$ for every $y$, while
$K(y,0)=f(y)$ is not identically zero because its first sine coefficient is
$a_1>0$.

\subsection{Exact SVE and failure of absolute convergence}

Let $T_K:L^2[-\pi,\pi]\to L^2[-\pi,\pi]$ be
$(T_K\varphi)(x)=\int_{-\pi}^{\pi}K(x,y)\varphi(y)\,dy$, and put
\[
 u_n(x)=\frac{\sin(nx)}{\sqrt\pi},
 \qquad
 v_n(y)=\frac{\cos(ny)}{\sqrt\pi}.
\]
Orthogonality and \eqref{eq:K} give
\begin{equation}\label{eq:svd}
 T_Kv_n=\pi a_nu_n,
 \qquad
 T_K^*u_n=\pi a_nv_n.
\end{equation}
Moreover, $T_K$ is Hilbert--Schmidt, vanishes on the constant and sine
subspaces, and its range is contained in the sine subspace.  Since the
constant, sine, and cosine systems give the standard orthogonal decomposition
of $L^2[-\pi,\pi]$, the complete list of positive
singular values is
\[
 \sigma_n=\pi a_n,
\]
and the SVE term $\sigma_nu_n(x)v_n(y)$ is exactly
$a_n\sin(nx)\cos(ny)$.  The sequence $(a_n)$ is strictly decreasing, so each
positive singular value is simple; only simultaneous sign changes of
$u_n,v_n$ are possible.

Finally, at $x=y=\pi/4$,
\begin{align*}
 \sum_{n=1}^{\infty}|\sigma_nu_n(x)v_n(y)|
 &=\sum_{n=1}^{\infty}a_n
   |\sin(n\pi/4)\cos(n\pi/4)|\\
 &=\frac12\sum_{\substack{n\geq1\\ n\ \mathrm{odd}}}
   \frac1{n\log(n+2)}=\infty.                                    \label{eq:diverge}
\end{align*}
For instance, after including the prefactor $1/2$, the $n=2k+1$ summand
dominates $1/[8(k+1)\log(k+2)]$, whose sum diverges.  This proves
Theorem~\ref{thm:main}.

\section{Verification, scope, and novelty}

Eight focused audits checked independently: positivity and $L^1$ convergence
of \eqref{eq:g}; every boundary term in \eqref{eq:mass}--\eqref{eq:coeff}; the
uniform tail estimate; identification of the primitive; both weak-variation
bounds; the exhaustive SVD; the simple-spectrum divergent point; and periodic
endpoint/scope issues.  The accompanying checker verifies high-precision
finite instances of the telescoping identities, monotonicity and convexity,
sampled uniform tails, simple coefficients, and growing absolute partial
sums.  These computations are sanity checks, not ingredients of the proof.

The counterexample lies in the paper's general asymmetric-kernel class and
therefore disproves the stated universal belief.  It does not address the
narrower question restricted to symmetric kernels.  It is compatible with
the authors' theorem: the failure occurs at one point, while their
unconditional conclusion is only almost everywhere.

Novelty is provisional.  Bounded searches through 13 August 2026 covered the
local run indexes and source corpus, the exact title and arXiv id, the quoted
uniform-BV/absolute-convergence statement, close Fourier/SVE phrases, the
arXiv and OpenReview records, and title/citation searches.  No later answer or
matching construction was found.  Priority is not claimed.  Human review
should first check the coefficient identity \eqref{eq:coeff} and the
exhaustiveness statement following \eqref{eq:svd}.

\begin{thebibliography}{9}
\bibitem{JeongTownsend}
S.~Jeong and A.~Townsend,
\emph{Extending Mercer's expansion to indefinite and asymmetric kernels},
Proceedings of ICLR 2025; arXiv:2409.16453v2.

\bibitem{Fejer}
L.~Fej\'er,
\emph{Untersuchungen \"uber Fouriersche Reihen},
Mathematische Annalen \textbf{58} (1904), 51--69.
\end{thebibliography}

\end{document}
